\subsection{Realizing the four-dimensional component obstruction}

We apply the realization theorem of Orson--Powell--Randal-Williams
to the component obstruction of
Theorem~\ref{thm:four-dimensional-component-image}.
Throughout this subsection $B$ is a connected, closed, oriented real
PL four-manifold, $N=B\times I$, and $\pi=\pi_1(B)$.
The spaces $A_{\mathrm{PL}}(B)$ and $P_{\mathrm{PL}}(B)$ use jointly
PL parameter families, while their TOP counterparts use the ordinary
compact-open topology. Concordances fix the bottom pointwise and have
free top. Write $c_A,c_P$ for the comparison maps and $e$ for top
restriction. Recall that
\[
 I_B=(\pi_0e)^{-1}\bigl(\operatorname{im}\pi_0c_A\bigr),
 \qquad
 \lambda_B:I_B\longrightarrow H^2(B;\mathbb Z/2)/J_B
\]
has kernel $\operatorname{im}\pi_0c_P$.
For a concordance whose two boundary maps are PL, its relative
structure difference is denoted by
\[
 \kappa(F)=[F_*\alpha_N]-[\alpha_N]
       \in H^3(N,\partial N;\mathbb Z/2)
       \cong H^2(B;\mathbb Z/2).
\]
Here one first imposes product collars and attaches exterior collars,
and uses the standard tangent lift as origin in the relative
open-five-dimensional PL-structure classification. The subgroup
$J_B$ consists of the classes $\kappa(T_\gamma)$ for TOP paths
$\gamma$ from the identity to a PL automorphism, where
$T_\gamma(x,t)=(\gamma_t(x),t)$.

Let $\delta_*$ be the Bockstein for
$0\to\mathbb Z_{(2)}\xrightarrow{2}\mathbb Z_{(2)}
\to\mathbb Z/2\to0$, and set
\begin{equation}\label{eq:four-dimensional-bockstein-subspace}
 L_B=\operatorname{PD}_B^{-1}
       \ker\bigl(\delta_*:H_2(B;\mathbb Z/2)
                      \longrightarrow H_1(B;\mathbb Z_{(2)})\bigr).
\end{equation}
We use precisely the quadratic assembly summand
\begin{equation}\label{eq:four-dimensional-assembly-summand}
 I_2:H_2(\pi;\mathbb Z_{(2)})
          \longrightarrow L^q_6(\mathbb Z[\pi])_{(2)}
\end{equation}
of \cite[Definition~5.8, p.~21]{OrsonPowellRandalWilliams2026}.
More explicitly, let $\mathcal L^q=\tau_{\geq0}L^q(\mathbb Z)$.
In the Taylor--Williams coordinates
\[
 H_6(B\pi;\mathcal L^q)_{(2)}\cong
 H_0(\pi;\mathbb Z/2)\oplus H_2(\pi;\mathbb Z_{(2)})
 \oplus H_4(\pi;\mathbb Z/2)\oplus H_6(\pi;\mathbb Z_{(2)}),
\]
the map is $I_2(b)=A_\pi\operatorname{TW}^{-1}(0,b,0,0)$,
where $A_\pi$ is quadratic assembly in the simple-surgery convention.

\begin{theorem}\label{thm:four-dimensional-obstruction-realization}
If $I_2=0$, then
\begin{equation}\label{eq:four-dimensional-obstruction-range}
       \operatorname{im}\lambda_B\supseteq(L_B+J_B)/J_B.
\end{equation}
If, in addition, $H_1(B;\mathbb Z_{(2)})$ has no $2$-torsion,
then $\lambda_B$ is surjective and induces an isomorphism
\begin{equation}\label{eq:four-dimensional-realized-quotient}
 I_B/\operatorname{im}(\pi_0c_P)
                 \cong H^2(B;\mathbb Z/2)/J_B.
\end{equation}
\end{theorem}

\begin{proof}
Choose a compatible smoothing of $B$, whose existence in dimension
four is recalled in \cite[Essay~I, p.~39]{KirbySiebenmann1977}.
Let $\sigma$ be its product smoothing on $N$, with underlying PL
structure $\alpha_N$.

\medskip\noindent\textit{Comparison of the relative classes.}
First let $F:N\to N$ fix both ends. Normalize it to be the identity
on common end collars. Put
\[
 W=N\times I,\qquad
 C=(N\times\{0\})\cup(\partial N\times I).
\]
The OPR boundary smoothing is $\sigma$ on $C$ and $F^*\sigma$ on
$N\times\{1\}$. Its relative Kirby--Siebenmann obstruction, followed
by Poincar\'e duality, is the invariant
$\operatorname{KS}_{\mathrm{OPR}}(F)\in H_2(B;\mathbb Z/2)$
\cite[pp.~2, 7--9]{OrsonPowellRandalWilliams2026}.

The universal square for $B\mathrm{PL}\to B\mathrm{TOP}$ identifies
a PL lift with a nullhomotopy of the universal map
$B\mathrm{TOP}\to K(\mathbb Z/2,4)$
\cite[Essay~V, Section~5.6, pp.~252--253]{KirbySiebenmann1977}.
With the standard lift fixed on $W$, changing the lift on the top
face by
\[
 d=[F^*\alpha_N]-[\alpha_N]
       \in H^3(N,\partial N;\mathbb Z/2)
       \cong H^3(\partial W,C;\mathbb Z/2)
\]
changes the relative obstruction by the connecting homomorphism
\begin{equation}\label{eq:relative-ks-interval-suspension}
 H^3(\partial W,C;\mathbb Z/2)
       \xrightarrow{\ \partial\ }
 H^4(W,\partial W;\mathbb Z/2).
\end{equation}
Indeed, the difference of the two boundary nullhomotopies is a
relative map to $K(\mathbb Z/2,3)$; its extension obstruction is
the connecting class in the cohomology sequence of the triple.
The inclusion $C\hookrightarrow W$ is a homotopy equivalence,
as simultaneous contraction in the last interval retracts both
spaces to $N\times\{0\}$. Thus
\eqref{eq:relative-ks-interval-suspension} is an isomorphism,
and is the interval suspension. Forgetting a smooth tangent lift
to a PL lift gives this same universal obstruction, so this
comparison identifies the classes themselves, not merely their
vanishing loci.

Since $F^*\alpha_N=(F^{-1})_*\alpha_N$, we have
$d=\kappa(F^{-1})$. A bottom-fixed concordance acts trivially on
the relative cohomology of $(N,\partial N)$. For completeness, if
$F(x,t)=(a(x,t),b(x,t))$, a homotopy of maps of pairs from $F$ to
the identity is
\[
 (x,t)\longmapsto
 \bigl(a(x,(1-s)t),(1-s)b(x,t)+st\bigr),\qquad 0\leq s\leq1.
\]
This is a homotopy of maps, and need not be an isotopy.
Additivity of the relative structure class therefore gives
$\kappa(F^{-1})=-\kappa(F)=\kappa(F)$.
Combining the two interval suspensions with mod-$2$ duality proves
\begin{equation}\label{eq:opr-relative-pl-identification}
 \operatorname{KS}_{\mathrm{OPR}}(F)
                    =\operatorname{PD}_B\kappa(F).
\end{equation}
The collar normalization is relative-boundary isotopic to the
original map and preserves both relative lift classes.

\medskip\noindent\textit{Realization with the ends fixed.}
Let $Q(B)$ be the OPR group of TOP pseudo-isotopies from the
identity to a diffeomorphism, with the equivalence relation in
\cite[Definition~2.1, p.~5]{OrsonPowellRandalWilliams2026}.
Theorem~2.3(D) of that paper says, under $I_2=0$, that the image of
\[
 \pi_1\widetilde{\operatorname{Homeo}}^+(B)
        \longrightarrow Q(B)
        \xrightarrow{\operatorname{KS}_{\mathrm{OPR}}}
                         H_2(B;\mathbb Z/2)
\]
contains $\ker\delta_*$.
The source is the block homeomorphism object. The assembly hypothesis
is the quadratic map in
\eqref{eq:four-dimensional-assembly-summand}, as specified in
\cite[Definition~5.8 and proof, pp.~21--22]{OrsonPowellRandalWilliams2026}.

Fix $z\in\ker\delta_*$. Choose a realizing element of $Q(B)$ and
represent it by a concordance $F$ from the identity to a
diffeomorphism $f$. Exactness of the block homotopy-fibre sequence
in \cite[equation~(1), p.~6]{OrsonPowellRandalWilliams2026}
implies that $f$ is smoothly pseudo-isotopic to the identity.
Choose such a smooth pseudo-isotopy $G$, with bottom $f$ and top
the identity. Then
\[
       G'=G\circ(f^{-1}\times\operatorname{id}_I),
       \qquad F_0=G'\circ F
\]
has bottom and top equal to the identity. The OPR invariant is a
homomorphism and vanishes on smooth concordances
\cite[Lemma~3.2 and the proof of Lemma~2.2, pp.~7--9]
{OrsonPowellRandalWilliams2026}. Consequently
\begin{equation}\label{eq:both-ended-obstruction-realization}
 \operatorname{KS}_{\mathrm{OPR}}(F_0)=z,
       \qquad \kappa(F_0)=\operatorname{PD}_B^{-1}(z).
\end{equation}
Smoothness of $G'$ is used to make its smoothing obstruction
zero. Its coordinate map need not be PL for the fixed
triangulation; the boundary maps of $F_0$ are literally the
identity, which is the condition needed in
\eqref{eq:opr-relative-pl-identification}.

The ordinary free-top component of $F_0$ lies in $I_B$, and its
image under $\lambda_B$ is
$\operatorname{PD}_B^{-1}(z)+J_B$. This proves
\eqref{eq:four-dimensional-obstruction-range}. If
$H_1(B;\mathbb Z_{(2)})$ has no $2$-torsion, exactness of the
coefficient sequence implies $\delta_*=0$. Hence
$L_B=H^2(B;\mathbb Z/2)$, and the previously established kernel
of $\lambda_B$ proves
\eqref{eq:four-dimensional-realized-quotient}.
\end{proof}

\begin{corollary}\label{cor:simply-connected-four-dimensional-trace}
If $B$ is a connected, closed, simply connected real PL
four-manifold, then
\[
              \pi_0P_{\mathrm{TOP}}(B)=1,
              \qquad J_B=H^2(B;\mathbb Z/2).
\]
For any such $B$, the ordinary trace homomorphism
\begin{equation}\label{eq:ordinary-four-dimensional-loop-obstruction}
 \eta_B:\pi_1A_{\mathrm{TOP}}(B)
       \longrightarrow H^2(B;\mathbb Z/2),
       \qquad [\gamma]\longmapsto\kappa(T_\gamma)
\end{equation}
is surjective and vanishes on the image of
$\pi_1A_{\mathrm{PL}}(B)$. Thus the ordinary automorphism
comparison is not surjective on $\pi_1$ if
$H^2(B;\mathbb Z/2)\ne0$.
\end{corollary}

\begin{proof}
Choose an orientation of $B$. Since $\pi=1$ and
$H_1(B;\mathbb Z_{(2)})=0$, both hypotheses of the theorem hold.
The corrected topological pseudo-isotopy theorem states that
every pseudo-isotopy of a compact simply connected four-manifold
is isotopic to the identity relative to the bottom and vertical
boundary \cite[Theorem~1.5, p.~3; conventions, p.~1]
{GabaiGayHartmanKrushkalPowell2026}. For closed $B$ this gives
$\pi_0P_{\mathrm{TOP}}(B)=1$.

Given $v\in H^2(B;\mathbb Z/2)$, choose the both-end-fixed
$F_0$ in \eqref{eq:both-ended-obstruction-realization} with
$\kappa(F_0)=v$. Take an ordinary path $H_s$ in
$P_{\mathrm{TOP}}(B)$ from the identity to $F_0$. The top maps
$\gamma_s=e(H_s)$ form a based ordinary loop, because both
ends of $F_0$ are the identity. Set
\[
 T_s(x,t)=(\gamma_{st}(x),t),\qquad Q_s=T_s^{-1}\circ H_s.
\]
The path $Q_s$ fixes both boundary components and joins the
identity to $T_\gamma^{-1}F_0$. Normalize the whole compact
parameter family on common end collars. Relative isotopy
invariance and additivity of $\kappa$ give
\[
            0=\kappa(T_\gamma^{-1}F_0)
                  =\kappa(T_\gamma)+\kappa(F_0).
\]
Hence $\eta_B[\gamma]=v$. This passage uses the ordinary
pseudo-isotopy theorem to obtain ordinary loops from the
realization statement.

Homotopies of based loops give relative-boundary isotopies of
their traces, so $\eta_B$ is well defined. Pointwise product of
loops induces product of traces; its operation on $\pi_1$ is
the usual loop operation. Additivity of $\kappa$ therefore
makes $\eta_B$ a homomorphism. A jointly PL loop has a jointly
PL trace, so its relative PL class is zero. Finally ordinary
loops are among the paths defining $J_B$, and surjectivity of
$\eta_B$ gives $J_B=H^2(B;\mathbb Z/2)$.
\end{proof}

\begin{example}\label{ex:four-dimensional-product-sphere-trace}
For $B=S^2\times S^2$, the target of
\eqref{eq:ordinary-four-dimensional-loop-obstruction} is
$(\mathbb Z/2)^2$. A finite PL model is the ordered product of
two tetrahedron boundaries. Its simplices are chains in the
coordinatewise order on $\{0,1,2,3\}^2$ whose two projections
have at most three vertices. These are the compatible ordered
subdivisions of the products of faces, so they triangulate the
product PL manifold itself.

One can count the $k$-simplices by the numbers $r,s$ of vertices
in their two projections. The successive steps increase the
first coordinate, the second coordinate, or both. Their numbers
are respectively $k+1-s$, $k+1-r$, and $r+s-k-2$. Thus
\[
 f_k=\sum_{r,s=1}^{3}\binom4r\binom4s
       \binom{k}{\,k+1-s,\ k+1-r,\ r+s-k-2\,},
 \qquad 0\leq k\leq4,
\]
where a term with a negative lower entry is zero. This gives
$f=(16,84,216,240,96)$. The product CW decomposition has one
$0$-cell, two $2$-cells and one $4$-cell, with zero cellular
differential; hence its integral homology is free of ranks
$(1,0,2,0,1)$. It is simply connected because each factor is.
The corollary therefore applies to this explicit finite base.
\end{example}

For general four-manifolds, the endpoint subgroup $I_B$, the trace
subgroup $J_B$, and the higher concordance groups remain to be determined.
